\begin{figure}[ht]
\centering
\begin{tikzpicture}[x=1.0cm, y=1.0cm, font=\small]

% Axes
\draw[->, thick] (0, 0) -- (10.6, 0)
  node[anchor=west, font=\scriptsize] {column-scale ratio $\log_{10}(s_{2}/s_{1})$};
\draw[->, thick] (0, 0) -- (0, 6.4)
  node[anchor=south, font=\scriptsize] {$\log_{10}\kappa$};

% x ticks
\foreach \x/\lab in {1/-4, 3/-2, 5/0, 7/2, 9/4} {
  \draw (\x, 0.08) -- (\x, -0.08) node[anchor=north, font=\scriptsize] {\lab};
}
% y ticks
\foreach \y in {0,2,4,6} {
  \draw (0.08, \y) -- (-0.08, \y) node[anchor=east, font=\scriptsize] {\y};
}

% U-shaped conditioning curve: minimum at balanced scaling (x=5), rising both ways
\draw[very thick, engblue]
  (1, 5.6) .. controls (3, 3.0) and (4.2, 1.1) ..
  (5, 0.7) .. controls (5.8, 1.1) and (7, 3.0) .. (9, 5.6);

% minimum marker
\fill[engred] (5, 0.7) circle (2pt);
\node[engred, anchor=south, font=\scriptsize] at (5, 0.85) {equilibrated: $\kappa$ minimal};

% unscaled point off to one side
\fill[engamber] (8, 4.4) circle (2pt);
\node[engamber, anchor=south west, font=\scriptsize] at (7.2, 4.5) {as-supplied units};
\draw[->, engamber, >=stealth, dashed] (7.9, 4.2) -- (5.3, 0.95);
\node[engamber, font=\scriptsize, anchor=west] at (5.6, 2.6) {rescale};

\node[anchor=north, font=\scriptsize, text width=11cm, align=center]
  at (5.3, -0.9) {
  Condition number against the relative scale of two matrix columns.
  A matrix whose columns carry mismatched physical units (a stiffness
  in $\si{\newton\per\meter}$ beside one in
  $\si{\newton\meter\per\radian}$) sits far up one arm; dividing each
  column by its norm moves it toward the basin, recovering digits the
  units had thrown away.
};

\end{tikzpicture}
\caption[Column scaling and the condition number]{Condition number of a
two-column matrix as a function of the ratio between the columns'
scales. Mismatched column magnitudes, common when the columns carry
different physical units, push the matrix up either arm of the curve and
inflate $\kappa$ for a reason that has nothing to do with the problem's
true sensitivity. Column equilibration, dividing each column by its norm
before the solve, moves the matrix toward the basin where $\kappa$ is
smallest, and the corresponding row or column scaling of the right-hand
side leaves the solution unchanged. Scaling is the cheapest defence
against artificial ill-conditioning, and a solver that reports $\kappa$
should report it after equilibration, not before.}
\label{fig:vol02:ch09:scaling-conditioning}
\end{figure}
